\documentclass{article}

\usepackage{amsmath} % math stuff
\usepackage{amssymb} % math stuff
\usepackage{array} % equations and stuff
\usepackage{bm} % bold math
%\usepackage{booktabs} % extra table rule options
%\usepackage{caption} % suppressed table numbering; incompatible with revtex, and longtable, I think
\usepackage{comment} % comment environment
%\usepackage{enumitem} % customization of enumeration, itemize, and description
\usepackage[T1]{fontenc} % font encoding for special characters, must also use scalable font package
\usepackage[margin=0.8in]{geometry} % paper sizes and margins (but be careful not to mess up pre-defined pages)
\usepackage{graphicx} % for graphics
%\usepackage{helvet} % default font is the helvetica postscript font
\usepackage[utf8]{inputenc} % special characters in tex input
\usepackage{layouts} % print units like widths
\usepackage{lipsum} % lorem ipsum filler text
\usepackage{lmodern} % scalable font?
\usepackage{longtable} % multi-page tables
\usepackage{makecell} % specify line-breaks in table cells
\usepackage{mathrsfs} % math script font
\usepackage{mhchem} % easier chemical formula
\usepackage{microtype} % allows disabling of ligatures
\usepackage{multicol} % multicolumns
%\usepackage{newcent} % new century schoolbook font
\usepackage{nicefrac}
\usepackage{numprint} % print and format (large) numbers
\usepackage{parskip} % removes paragraph indentation, and adjusts paragraph skip, as well as list items
\usepackage{pdfpages} % add pdf files as pages
%\usepackage{setspace} % adjust text spacing and indents
\usepackage{siunitx} % decimal alignment, and degree symbol
\usepackage{subfigure} % divided figures
%\usepackage{tabu} % extra table options
\usepackage{textcomp} % symbols
\usepackage{threeparttablex} % better footnotes with longtable
\usepackage{titling} % title placement
\usepackage{ulem} % strikethrough text
\usepackage{upgreek} % upright Greek letters
%\usepackage{url} % superceded by hyperref
\usepackage{verbatim} % verbatim environment
\usepackage{xcolor} % colors and color boxes
\usepackage{xspace} % commands that don't eat up white space
\usepackage{hyperref} % links and page setup; should always come last

\hypersetup{
 bookmarks=true,
 colorlinks=true,
 citecolor=blue,
 linkcolor=blue,
 urlcolor=blue,
 pdfstartview={XYZ null null 1.0} % default open view is 100%
}

\DisableLigatures[f,t]{encoding = T1} % disable ff, fi, fl, tt ligatures; without options, it also disables -- = endash
\renewcommand{\arraystretch}{1.0} % extra vertical (and horizontal?) space in tables

% define centered, left- and right-aligned columns with specified widths
\newcommand{\PreserveBackslash}[1]{\let\temp=\\#1\let\\=\temp}
\newcolumntype{C}[1]{>{\PreserveBackslash\centering}p{#1}}
\newcolumntype{L}[1]{>{\PreserveBackslash\raggedright}p{#1}}
\newcolumntype{R}[1]{>{\PreserveBackslash\raggedleft}p{#1}}

\begin{document}

\pagestyle{empty} % don't number pages

% custom title
\begin{center}
{\LARGE Classic Riddler}

\vspace{0.15in}

{\Large 15 April 2022}
\end{center}


\section*{Riddle:}

For this riddle, we'll define $N$ as being a Carmichael number if (1) it is composite and has no square factors, and (2) if one less than every prime factor of $N$ is a factor of $N-1$.

That's a lot to take in at once, so let's look at an example.
The smallest Carmichael number is 561. Sure enough, it has no square factors (other than 1, which we're not counting here).
Meanwhile, its prime factors are 3, 11 and 17.
If we subtract one from each of those, we get 2, 10 and 16, all of which divide 560 (one less than the original number).

Daniel's puzzle asks you to find a six-digit Carmichael number that can be written as ABCABC in base 10.
(\textbf{Note:} The digits A, B and C are \textit{not} necessarily distinct.)
Of course, you could look this up in a matter of seconds or even write some code to figure it out.
But I assure you, this riddle can be done by hand, so please give it a shot!


\section*{Solution:}

The number ABCABC can easily be factored into $1001\times\mathrm{ABC}$.
1001 itself can be further factored into $7\times11\times13$.
So at least 7, 11, and 13 are factors, and none of 7, 11, or 13 can divide ABC.
In order to be a Carmichael number, 6, 10, and 12 must divide $\mathrm{ABCABC}-1$.
Because $\mathrm{ABCABC}-1$ is divisible by 10, $\mathrm{C}-1=0$, so C must be 1.
Because AB1AB0 is divisible by 3 (because it is divisible by 6), neither AB1AB1 nor AB1 is divisible by 3.

There are 60 combinations of A and B where 3 is not a factor of AB1 (101, 121, 131, 151, etc.).
This is few enough that they can be checked by hand.
Luckily, the lowest combination, 101 meets the Carmichael criteria.
101 itself is prime, so the prime factorization of 101,101 is $7\times11\times13\times101$.
101,100 must therefore be divisible by 6, 10, 12, and 100.
It trivially is divisible by 10 and 100 $(=2\times2\times5\times5)$, as well as 3.
Because of the factor of 3 and two factors of 2, it must also be divisible by 6 and 12.

So the value of \fcolorbox{red}{white}{\textbf{101,101}} is a solution to the riddle.
I don't know if there are additional solutions.



\end{document}